% SPDX-License-Identifier: MIT
% Copyright (c) 2026 Evgeny Baulin
% Optimization Methods, Seminar 10: theory handout (Russian).
% Generated by checks/10. Knapsack and Linear Programming in Practice.ipynb from web/10/js/content-*.js and content/students.json. Do not edit by hand.
% Compile with: latexmk -xelatex -interaction=nonstopmode -halt-on-error Theory_ru.tex
\documentclass[11pt,a4paper]{article}

\usepackage{fontspec}
\setmainfont{cmunrm.otf}[BoldFont=cmunbx.otf, ItalicFont=cmunti.otf, BoldItalicFont=cmunbi.otf]
\setsansfont{cmunss.otf}[BoldFont=cmunsx.otf, ItalicFont=cmunsi.otf, BoldItalicFont=cmunso.otf]
\setmonofont{cmuntt.otf}[BoldFont=cmuntb.otf, ItalicFont=cmunit.otf, BoldItalicFont=cmuntx.otf]
\usepackage{polyglossia}
\setdefaultlanguage{russian}
\setotherlanguage{english}
\newfontfamily\cyrillicfont{cmunrm.otf}[Ligatures=TeX, BoldFont=cmunbx.otf, ItalicFont=cmunti.otf, BoldItalicFont=cmunbi.otf]
\newfontfamily\cyrillicfontsf{cmunss.otf}[Ligatures=TeX, BoldFont=cmunsx.otf, ItalicFont=cmunsi.otf, BoldItalicFont=cmunso.otf]
\newfontfamily\cyrillicfonttt{cmuntt.otf}[BoldFont=cmuntb.otf, ItalicFont=cmunit.otf, BoldItalicFont=cmuntx.otf]

\usepackage[a4paper,left=2.3cm,right=2.3cm,top=2.4cm,bottom=2.4cm,headheight=14pt]{geometry}
\usepackage{amsmath,amssymb}
\usepackage{graphicx}
\usepackage{array,booktabs,longtable}
\usepackage{enumitem}
\usepackage{xcolor}
\usepackage[most]{tcolorbox}
\usepackage{titlesec}
\usepackage{fancyhdr}
\usepackage{needspace}
\usepackage[expansion=false]{microtype}
\usepackage[colorlinks=true,linkcolor=hseblue,urlcolor=hseblue,citecolor=hseblue]{hyperref}

\definecolor{hseblue}{RGB}{16,45,105}
\definecolor{hseaccent}{RGB}{198,40,40}
\definecolor{hsegray}{RGB}{108,114,126}
\definecolor{hselight}{RGB}{240,243,248}
\definecolor{hsewarm}{RGB}{253,244,232}

\hypersetup{pdftitle={Методы оптимизации. Семинар 10. Конспект теории}, pdfauthor={Евгений Баулин}}

\titleformat{\section}{\color{hseblue}\Large\bfseries}{\thesection}{0.6em}{}
\titleformat{\subsection}{\color{hseblue}\large\bfseries}{\thesubsection}{0.6em}{}
\titlespacing*{\section}{0pt}{16pt plus 4pt minus 2pt}{8pt plus 2pt}
\titlespacing*{\subsection}{0pt}{12pt plus 3pt minus 2pt}{6pt plus 2pt}
\pretocmd{\section}{\Needspace{7\baselineskip}}{}{\errmessage{cannot patch section}}
\pretocmd{\subsection}{\Needspace{5\baselineskip}}{}{\errmessage{cannot patch subsection}}

\tcbset{handoutbox/.style={enhanced, breakable, sharp corners, boxrule=0pt, left=7pt, right=7pt, top=5pt, bottom=5pt, before skip=8pt plus 2pt, after skip=8pt plus 2pt}}
\newtcolorbox{examplebox}{handoutbox, colback=hselight, colframe=hselight, before upper={\textcolor{hseblue}{\bfseries Пример.}\ }}
\newtcolorbox{whybox}{handoutbox, colback=white, colframe=hseblue, leftrule=2.5pt, top=3pt, bottom=3pt, before upper={\textcolor{hseaccent}{\bfseries Почему.}\ }}
\newtcolorbox{answerbox}{handoutbox, colback=hsewarm, colframe=hsewarm, before upper={\textcolor{hseaccent}{\bfseries Ответ.}\ }}
\newtcolorbox{conventionbox}{handoutbox, colback=hselight, colframe=hseblue, leftrule=2.5pt, before upper={\textcolor{hseblue}{\bfseries Обозначения.}\par\smallskip}}

\newcommand{\handoutfigure}[3]{\par\medskip\noindent\begin{minipage}{\linewidth}\centering
  \includegraphics[width=#1\linewidth]{figures/#2.pdf}\par\smallskip
  \begin{minipage}{0.94\linewidth}\small\textcolor{hseblue}{\bfseries Рисунок.}\ #3\end{minipage}\end{minipage}\par\medskip}
\newcommand{\problemhead}[3]{\par\bigskip\noindent{\color{hseblue}\large\bfseries #1.\ #2}\par\nopagebreak
  \noindent{\small\itshape\color{hsegray}#3}\par\nopagebreak\smallskip}
\newtcolorbox{definitionbox}[1]{handoutbox, colback=white, colframe=hseblue, leftrule=2.5pt, before upper={\textcolor{hseblue}{\bfseries #1}\ }}
\newtcolorbox{theorembox}[1]{handoutbox, colback=hselight, colframe=hseblue, leftrule=2.5pt, before upper={\textcolor{hseblue}{\bfseries #1}\ }}
\newtcolorbox{proofbox}{handoutbox, colback=white, colframe=hsegray, leftrule=1.2pt, top=3pt, bottom=3pt, before upper={\textcolor{hsegray}{\bfseries\itshape Доказательство.}\ }}
\newtcolorbox{notebox}[1]{handoutbox, colback=hselight, colframe=hselight, before upper={\textcolor{hseblue}{\bfseries #1}\par\smallskip}}
\newtcolorbox{trapbox}{handoutbox, colback=hsewarm, colframe=hseaccent, leftrule=2.5pt, before upper={\textcolor{hseaccent}{\bfseries Ловушка.}\ }}
\newtcolorbox{selfcheckbox}{handoutbox, colback=white, colframe=hseaccent, leftrule=2.5pt, top=3pt, bottom=3pt, before upper={\textcolor{hseaccent}{\bfseries Проверьте себя.}\ }}
\newcommand{\handoutdisplay}[1]{\sbox0{$\displaystyle #1$}\ifdim\wd0>\linewidth
  \par{\centering\binoppenalty=10000\relpenalty=0\relax$\displaystyle #1$\par}\else\[#1\]\fi}
\newcommand{\selfcheckanswer}{\par\smallskip\textcolor{hsegray}{\bfseries Ответ.}\ }

\setlist[itemize]{itemsep=3pt,topsep=4pt}
\setlist[enumerate]{itemsep=3pt,topsep=4pt}
\setlength{\parindent}{0pt}
\setlength{\parskip}{5pt plus 1pt}
\setlength{\emergencystretch}{2em}
\setcounter{tocdepth}{2}

\pagestyle{fancy}
\fancyhf{}
\fancyhead[L]{\small\color{hsegray}Методы оптимизации}
\fancyhead[R]{\small\color{hsegray}Семинар 10}
\fancyfoot[L]{\small\color{hsegray}Евгений Баулин}
\fancyfoot[C]{\small\thepage}
\fancyfoot[R]{\small\color{hsegray}Лицензия MIT}
\renewcommand{\headrulewidth}{0pt}
\renewcommand{\footrulewidth}{0pt}
\fancypagestyle{plain}{\pagestyle{fancy}}

\begin{document}

\begin{center}
  {\large\color{hsegray}Методы оптимизации\par}
  \vspace{4pt}
  {\large\bfseries\color{hseaccent}Семинар 10\par}
  \vspace{8pt}
  {\LARGE\bfseries\color{hseblue}Рюкзак, оценки и целочисленные модели\par}
  \vspace{8pt}
  {\Large Конспект теории\par}
  \vspace{12pt}
  НИУ ВШЭ, факультет компьютерных наук\par
  2026-2027 учебный год\par
  \vspace{8pt}
  Лектор курса: Андрей Игнатов\par
  Евгений Баулин\par
  Telegram: \href{https://t.me/tarakan_tuc}{@tarakan\_tuc}\par
\end{center}

\noindent{\color{hseblue}\rule{\linewidth}{0.6pt}}

\tableofcontents

\clearpage


Этот конспект сопровождает семинар 10. Он начинается с рюкзака 0–1 и его релаксации, строит на нём метод ветвей и границ и динамическое программирование, а затем переходит к моделированию двоичными переменными. Все задачи семинара приведены здесь с полными решениями.

\begin{conventionbox}
Максимизация, если не сказано иное. Ценности $p_i>0$, веса $w_i>0$, вместимость $W$; $x_i\in\{0,1\}$ говорит, уложен ли предмет $i$. $U(N)$ — верхняя оценка для узла $N$ перебора, $z_{\text{inc}}$ — значение лучшей найденной допустимой точки. $D(i,c)$ — лучшая ценность, которую дают первые $i$ предметов, если вместимость равна $c$.
\end{conventionbox}

% =====================================================================
\section{Рюкзак и оценка по плотности}\label{sec:knapsack}

Рюкзак — простейшая целочисленная задача с одним ограничением, и на нём видно почти всё о целочисленной оптимизации: оценки, перебор, таблицы и разрыв между релаксацией и целыми.

\begin{definitionbox}{Определение: Рюкзак 0–1 и его релаксация.}
$\max\sum_ip_ix_i$ при $\sum_iw_ix_i\le W$, $x\in\{0,1\}^n$. Релаксация сохраняет цель и строку и заменяет $x_i\in\{0,1\}$ на $0\le x_i\le1$. Варианты: ограниченные копии $0\le x_i\le u_i$ целые, неограниченные копии $x_i\in\mathbb Z_{\ge0}$, несколько строк ресурсов, несколько рюкзаков.
\end{definitionbox}

\begin{theorembox}{Теорема: Дробный рюкзак.}
Упорядочим предметы по убыванию плотности $p_i/w_i$. Берём целые предметы, пока они помещаются, затем долю следующего, заполняющую вместимость. Это решение релаксации, и дробной в нём будет не более одной переменной.
\end{theorembox}

\begin{proofbox}
Возьмём оптимальное решение релаксации, в котором более плотный предмет $i$ взят не полностью, а менее плотный $j$ занимает вес. Перенос малого веса $\delta$ от $j$ к $i$ сохраняет строку и меняет ценность на $\delta(p_i/w_i-p_j/w_j)\ge0$. Повторяя это, придём к полностью взятому началу упорядоченного списка и одному частичному предмету — это и есть жадное решение.
\end{proofbox}

\begin{whybox}
Релаксация — это продажа на вес: платим вместимостью, поэтому сначала берём товар с лучшей ценой за килограмм. Целочисленность всё ломает: последний предмет нельзя разрезать.
\end{whybox}

\handoutfigure{0.72}{fig_a1}{Дрон из A1: вместимость, занятая предметами по убыванию плотности, и под ней оптимальная загрузка.}

\problemhead{A1}{Посылки для дрона}{У доски}
\textbf{Условие.} Пять посылок приносят за доставку $(17;\ 12;\ 7;\ 19;\ 17)$ и весят $(8;\ 6;\ 7;\ 9;\ 4)$ кг; дрон поднимает не больше $W=19$ кг. Решите релаксацию по плотностям, округлите дробную посылку вниз и найдите лучшую загрузку. Укажите оценку $U$, ценность округлённой загрузки и оптимум.

\textbf{Решение.}
\begingroup\sloppy
\begin{enumerate}[label=\textbf{\arabic*.}, leftmargin=*]
  \item Плотности равны $(\tfrac{17}{8};\ 2;\ 1;\ \tfrac{19}{9};\ \tfrac{17}{4})$, поэтому порядок такой: $(5;\ 1;\ 4;\ 2;\ 3)$. Берём посылки целиком в этом порядке, пока очередная не перестанет помещаться.
  \item Посылки 5 и 1 помещаются, посылка 4 — нет, и оставшаяся вместимость берёт её долю $\tfrac{7}{9}$. Решение релаксации: $x=(1,0,0,\tfrac{7}{9},1)$, $U=\tfrac{439}{9}\approx48{,}7778$.
  \item Если округлить $x_{ 4 }$ вниз, останется загрузка ценностью $34$ и много свободного места. Если пропустить непоместившуюся посылку и идти по списку дальше, получим $46$ с посылками $(5;\ 1;\ 2)$: лучше, но всё ещё не оптимум.
  \item Посылки $(2;\ 4;\ 5)$ весят ровно $W$ и приносят $48$. Ценности целые, поэтому никакая загрузка не приносит больше $\lfloor U\rfloor$, а это тоже $48$. Оценка доказывает, что эта загрузка оптимальна.
\end{enumerate}
\endgroup
\begin{answerbox}
$U=\tfrac{439}{9}\approx48{,}7778$, округлённая загрузка $34$, оптимум $48$ с посылками $(2;\ 4;\ 5)$.
\end{answerbox}

\begin{trapbox}
Округление решения релаксации вниз — эвристика, а не метод: в A1 оно теряет почти треть ценности. При целых ценностях честная верхняя оценка — $\lfloor U\rfloor$.
\end{trapbox}

\problemhead{A2}{Вес и объём}{Самостоятельно}
\textbf{Условие.} Пять коробок приносят $(13;\ 7;\ 11;\ 16;\ 9)$, весят $(7;\ 6;\ 5;\ 9;\ 3)$ и занимают объёмы $(6;\ 2;\ 6;\ 6;\ 3)$; фургон везёт вес $W=17$ и объём $V=10$. Берите коробки в порядке ценности на единицу веса, пока позволяют оба предела, затем найдите лучшую загрузку и значение релаксации с обеими строками. Укажите значение жадного правила, оптимум и оценку ЛП.

\textbf{Решение.}
\begingroup\sloppy
\begin{enumerate}[label=\textbf{\arabic*.}, leftmargin=*]
  \item По ценности на единицу веса порядок такой: $(5;\ 3;\ 1;\ 4;\ 2)$. Идём по нему, пропуская то, что нарушает предел, и получаем жадную загрузку $(5;\ 3)$ ценностью $20$: лёгкие плотные коробки съедают объём.
  \item Перебирая укладки, проходящие по обоим пределам, находим лучшую: $(4;\ 5)$ ценностью $25$.
  \item Оптимум релаксации с обеими строками — в $x=(0,1,0,\tfrac{5}{6},1)$ со значением $\tfrac{88}{3}\approx29{,}3333$. Строка объёма там выполняется как равенство, а у строки веса остаётся запас, поэтому дробная переменная одна; если бы обе строки были равенствами, в вершине их могло бы быть две.
\end{enumerate}
\endgroup
\begin{answerbox}
Жадное правило $20$, оптимум $25$ с коробками $(4;\ 5)$, оценка ЛП $\tfrac{88}{3}\approx29{,}3333$.
\end{answerbox}

\problemhead{A3}{Копии одного предмета}{Самостоятельно}
\textbf{Условие.} Предметы трёх видов приносят $(7;\ 14;\ 12)$ и весят $(7;\ 3;\ 2)$, каждого вида не больше $u=4$ копий; вместимость $W=13$. Решите релаксацию $0\le x_i\le u$, затем найдите целый оптимум. Укажите оценку ЛП, оптимум и число копий второго вида в нём.

\textbf{Решение.}
\begingroup\sloppy
\begin{enumerate}[label=\textbf{\arabic*.}, leftmargin=*]
  \item Самый плотный — третий вид, за ним второй. Релаксация берёт $u$ копий третьего и такую долю второго, что заполняет остаток: $U=\tfrac{214}{3}\approx71{,}3333$.
  \item Если округлить это вниз и дозаполнить жадно целыми копиями, получится $62$. Если взять на две копии самого плотного вида меньше, освободится место ещё для двух копий второго: оптимум $x=(0;\ 3;\ 2)$ со значением $66$, и вместимость использована до последней единицы.
\end{enumerate}
\endgroup
\begin{answerbox}
$U=\tfrac{214}{3}\approx71{,}3333$, оптимум $66$ при $x=(0;\ 3;\ 2)$.
\end{answerbox}

\begin{selfcheckbox}
Почему сортировка по плотности не решает релаксацию с двумя строками ресурсов?
\selfcheckanswer Единой цены за единицу нет: предмет может быть лёгким, но объёмным. Релаксация — настоящая задача ЛП, и в её вершине могут быть две дробные переменные.
\end{selfcheckbox}

% =====================================================================
\section{Метод ветвей и границ}\label{sec:bnb}

Перебор всех $2^n$ укладок — двоичное дерево: каждый уровень фиксирует одну переменную. Метод ветвей и границ ходит по этому дереву, но пропускает всякое поддерево, о котором доказано, что в нём нет ничего лучше уже найденного.

\begin{definitionbox}{Определение: Узел, оценка, рекорд.}
Узел $N$ закрепляет значения некоторых переменных; его продолжения — укладки, согласные с ними. $U(N)$ — верная верхняя оценка, если у каждого допустимого продолжения значение не больше $U(N)$; для рюкзака это зафиксированная ценность плюс релаксация по свободным предметам в оставшейся вместимости. Рекорд $z_{\text{inc}}$ — лучшее найденное допустимое значение.
\end{definitionbox}

\begin{theorembox}{Теорема: Отсечение и точность.}
Если $U(N)\le z_{\text{inc}}$, в узле $N$ нет лучшего решения. При верных оценках и полном ветвлении рекорд оптимален, когда не осталось открытых узлов.
\end{theorembox}

\begin{proofbox}
У любого допустимого продолжения $x$ узла $N$ выполнено $f(x)\le U(N)\le z_{\text{inc}}$, поэтому, отбросив $N$, мы ничего не теряем. Каждое ветвление фиксирует ещё одну переменную, поэтому дерево конечно, и каждая допустимая точка лежит ровно в одном листе, который либо рассмотрен, либо отброшен по этой причине.
\end{proofbox}

\begin{itemize}
  \item Узел завершается в трёх случаях: закреплённые предметы превышают вместимость; $U(N)\le z_{\text{inc}}$; оптимум релаксации в узле уже является укладкой, и она может заменить рекорд.
  \item Ветвитесь по дробному предмету релаксации: оба потомка меняют оценку, и один из них часто сразу закрывается.
  \item В глубину: мало памяти и скорый рекорд. По лучшей оценке: меньше узлов до доказательства оптимальности, больше открытых узлов в памяти.
  \item Даже незаконченный прогон кое-что гарантирует: рекорд ограничивает оптимум снизу, а самая большая оценка среди открытых узлов — сверху.
\end{itemize}

\handoutfigure{0.66}{fig_b1}{Дерево перебора B1 со всеми оценками: крестик — недопустимый узел, оранжевый закрыт по оценке, красные — рекорды.}

\problemhead{B1}{Дерево перебора вручную}{У доски}
\textbf{Условие.} Ценности $(9;\ 14;\ 15;\ 14)$, веса $(6;\ 3;\ 4;\ 9)$, вместимость $W=13$. Проведите метод ветвей и границ в глубину. В каждом узле вычисляйте оценку по плотности, ветвитесь по дробному предмету, сначала по ветке $x_j=1$, и закрывайте узел, если он недопустим, целый или $U(N)\le z_{\text{inc}}$. Считайте каждый узел, для которого вычислена оценка, включая корень и недопустимые. Укажите число узлов и оптимум.

\textbf{Решение.}
\begingroup\sloppy
\begin{enumerate}[label=\textbf{\arabic*.}, leftmargin=*]
  \item Корень: предметы 2 и 3 помещаются, предмет 4 — лишь частично, поэтому $U=\tfrac{115}{3}\approx38{,}3333$, и мы ветвимся по $x_{ 4 }$.
  \item При $x_{ 4 }=1$ оценка равна $31{,}75$, а дробным становится предмет 3. Ветка $x_{ 3 }=1$ заполняет рюкзак ровно и даёт целую релаксацию: первый рекорд $z_{\text{inc}}=29$.
  \item Ветка $x_{ 3 }=0$ имеет $U=29{,}5$, выше рекорда, поэтому ветвимся по предмету 1. Если взять его, рюкзак переполнен (недопустимо); если нет, $U=28\le z_{\text{inc}}$, и узел закрыт по оценке.
  \item Возвращаемся к корню, $x_{ 4 }=0$: три остальных предмета заполняют рюкзак ровно, релаксация целая со значением $38$ — новый рекорд. Открытых узлов нет, значит, он оптимален. Весь журнал:
\handoutdisplay{
\begin{array}{c|l|c|c} k & \text{fixed} & U & z_{\text{inc}} \\ \hline 1 & \text{---} & 38{,}33 & 0 \\ 2 & x_4=1 & 31{,}75 & 0 \\ 3 & x_4=1;\ x_3=1 & 29 & 0 \\ 4 & x_4=1;\ x_3=0 & 29{,}5 & 29 \\ 5 & x_4=1;\ x_3=0;\ x_1=1 & \text{---} & 29 \\ 6 & x_4=1;\ x_3=0;\ x_1=0 & 28 & 29 \\ 7 & x_4=0 & 38 & 29 \end{array}
}

\end{enumerate}
\endgroup
\begin{answerbox}
Узлов 7, оптимум $38$ с предметами $(1;\ 2;\ 3)$.
\end{answerbox}

\problemhead{B2}{В глубину или по лучшей оценке}{Самостоятельно}
\textbf{Условие.} Ценности $(6;\ 7;\ 17;\ 19;\ 20)$, веса $(7;\ 4;\ 6;\ 5;\ 9)$, вместимость $W=18$. Проведите метод ветвей и границ по правилам B1 дважды: в глубину и по лучшей оценке, когда всегда раскрывается открытый узел с наибольшим $U$; там потомок закрывается сразу после вычисления оценки, если он недопустим, целый или не выше рекорда. Считайте узлы, для которых вычислена оценка, включая корень. Укажите оба числа и оптимум.

\textbf{Решение.}
\begingroup\sloppy
\begin{enumerate}[label=\textbf{\arabic*.}, leftmargin=*]
  \item Порядок по плотности: $(4;\ 3;\ 5;\ 2;\ 1)$, оценка корня $\tfrac{464}{9}$. Поиск в глубину идёт вниз по веткам $x_j=1$ и встречает три недопустимых потомка до первой укладки, которая даёт первый рекорд в журнале:
\handoutdisplay{
\begin{array}{c|l|c|c} k & \text{fixed} & U & z_{\text{inc}} \\ \hline 1 & \text{---} & 51{,}56 & 0 \\ 2 & x_5=1 & 50{,}33 & 0 \\ 3 & x_5=1;\ x_3=1 & 48{,}4 & 0 \\ 4 & x_5=1;\ x_3=1;\ x_4=1 & \text{---} & 0 \\ 5 & x_5=1;\ x_3=1;\ x_4=0 & 42{,}25 & 0 \\ 6 & x_5=1;\ x_3=1;\ x_4=0;\ x_2=1 & \text{---} & 0 \\ 7 & x_5=1;\ x_3=1;\ x_4=0;\ x_2=0 & 39{,}57 & 0 \\ 8 & x_5=1;\ x_3=1;\ x_4=0;\ x_2=0;\ x_1=1 & \text{---} & 0 \\ 9 & x_5=1;\ x_3=1;\ x_4=0;\ x_2=0;\ x_1=0 & 37 & 0 \\ 10 & x_5=1;\ x_3=0 & 46 & 37 \\ 11 & x_5=0 & 45{,}57 & 46 \end{array}
}
 Всего 11 узлов.
  \item Поиск по лучшей оценке раскрывает потомка корня с большей оценкой, и один из его потомков уже оказывается укладкой ценностью $46$. Затем он раскрывает единственный открытый узел, оценка которого ещё выше этого значения; его потомки недопустимы или ниже, а второй потомок корня, лежащий ниже рекорда, так и не раскрывается. Он вычисляет 7 оценок.
  \item Оба приходят к оптимуму $46$, предметы $(2;\ 4;\ 5)$.
\end{enumerate}
\endgroup
\begin{answerbox}
В глубину 11 узлов, по лучшей оценке 7, оптимум $46$.
\end{answerbox}

\begin{trapbox}
Для подсчёта узлов нужна договорённость: считаются ли корень, недопустимые потомки и сразу закрытые потомки. Сравнивайте стратегии только при одной и той же договорённости.
\end{trapbox}

% =====================================================================
\section{Отсекающие плоскости}\label{sec:cuts}

Вместо того чтобы делить задачу, её можно усилить: добавить неравенство, верное для каждой целой точки и нарушенное в найденном оптимуме релаксации. Солверы делают и то и другое, это называется методом ветвей и отсечений.

\begin{theorembox}{Теорема: Дробное отсечение Гомори.}
Пусть данные целые, а строка оптимальной таблицы имеет вид $x_k+\sum_{j\in N}\bar a_jx_j=\bar b$ с дробным $\bar b$, где $N$ — небазисные переменные. Тогда $\sum_{j\in N}\{\bar a_j\}x_j\ge\{\bar b\}$ верно во всех целых точках и нарушено в оптимуме ЛП, где $\{t\}=t-\lfloor t\rfloor$.
\end{theorembox}

\begin{proofbox}
Разделим каждый коэффициент на целую и дробную части: $x_k+\sum\lfloor\bar a_j\rfloor x_j-\lfloor\bar b\rfloor=\{\bar b\}-\sum\{\bar a_j\}x_j$. В целой точке левая часть целая, а правая меньше единицы, значит, она не больше нуля. В оптимуме ЛП все $x_j$, $j\in N$, равны нулю, и правая часть равна $\{\bar b\}>0$.
\end{proofbox}

\begin{whybox}
Отсечение срезает угол многогранника, в котором лежит дробный оптимум, и ничего важного больше. В двух переменных это видно: прямая $x_1+x_2\le d$ из B3 проходит через целые точки.
\end{whybox}

\problemhead{B3}{Одно отсечение Гомори}{Самостоятельно}
\textbf{Условие.} Максимизируйте $4x_1+2x_2$ при $6x_1+2x_2\le14$, $x_1+2x_2\le10$, $x\ge0$ целых. Решите релаксацию, возьмите строку $x_1$ оптимальной таблицы и выведите отсечение Гомори; в переменных $x$ оно имеет вид $x_1+x_2\le d$. Укажите значение ЛП, $d$ и значение ЛП после добавления отсечения.

\textbf{Решение.}
\begingroup\sloppy
\begin{enumerate}[label=\textbf{\arabic*.}, leftmargin=*]
  \item Обе строки как равенства дают $x=(\tfrac{4}{5},\tfrac{23}{5})$ и $z=\tfrac{62}{5}=12{,}4$; обе переменные дробные.
  \item Когда $x_1,x_2$ базисные, строка $x_1$ такова:
\handoutdisplay{
x_1+\tfrac{1}{5}s_1-\tfrac{1}{5}s_2=\tfrac{4}{5}
}
 Дробные части $-\tfrac{1}{5}$ и $\tfrac{1}{5}$ равны $\tfrac{4}{5}$ и $\tfrac{1}{5}$, поэтому отсечение:
\handoutdisplay{
\tfrac{1}{5}s_1+\tfrac{4}{5}s_2\ge\tfrac{4}{5}
}

  \item Подставим $s_1=b_1-6x_1-2x_2$ и $s_2=b_2-x_1-2x_2$: слагаемые с $x$ дают $-2x_1-2x_2$, и неравенство превращается в $x_1+x_2\le5$. У текущего оптимума $x_1+x_2>5$, поэтому он отсекается, а все целые точки отсечению удовлетворяют.
  \item С отсечением оптимум ЛП переходит в $(1;\ 4)$ со значением $12$. Эта точка целая, значит, это целый оптимум: одно отсечение закрыло разрыв.
\end{enumerate}
\endgroup
\begin{answerbox}
$z=\tfrac{62}{5}=12{,}4$, отсечение $x_1+x_2\le5$, после него $z=12$ в $(1;\ 4)$.
\end{answerbox}

% =====================================================================
\section{Динамическое программирование для рюкзака}\label{sec:dp}

При целых весах у рюкзака есть второй точный метод. Он не ищет по дереву, а заполняет таблицу подзадач, каждая из которых решается один раз из двух меньших.

\begin{definitionbox}{Определение: Состояние.}
$D(i,c)$ — лучшая ценность, если использовать только первые $i$ предметов при целой вместимости $c$, для $0\le i\le n$ и $0\le c\le W$. Граничные значения $D(0,c)=0$, ответ — $D(n,W)$.
\end{definitionbox}

\begin{theorembox}{Теорема: Рекуррентное соотношение.}
$D(i,c)=D(i-1,c)$, если $w_i>c$, и $D(i,c)=\max\{D(i-1,c),\ p_i+D(i-1,c-w_i)\}$ иначе.
\end{theorembox}

\begin{proofbox}
Оптимальный выбор среди первых $i$ предметов либо не берёт предмет $i$ и тогда оптимален для $(i-1,c)$, либо берёт его, и тогда остальное оптимально для $(i-1,c-w_i)$. Оба варианта достижимы, поэтому оптимален больший из них. Индукция по $i$ завершает доказательство.
\end{proofbox}

\begin{itemize}
  \item Обратный ход: начинаем с $(n,W)$; если $D(i,c)=D(i-1,c)$, пропускаем предмет $i$, иначе берём его и вычитаем $w_i$ из $c$.
  \item Для значения хватает одной строки памяти: обновляйте её от $c=W$ вниз до $c=w_i$, чтобы каждый предмет использовался один раз.
  \item Неограниченные копии: $D(i,c)=\max\{D(i-1,c),\ p_i+D(i,c-w_i)\}$, та же строка; в одной строке проходите вверх.
  \item Ограниченные копии: максимум по $q=0,\dots,\min(u_i,\lfloor c/w_i\rfloor)$ копиям предмета $i$.
\end{itemize}

\handoutfigure{0.72}{fig_c1}{Задача C1 таблицей; клетки в рамках прослеживают укладку назад.}

\problemhead{C1}{Таблица маленького рюкзака}{У доски}
\textbf{Условие.} Ценности $(13;\ 10;\ 6;\ 5)$, веса $(5;\ 4;\ 3;\ 2)$, вместимость $W=8$. Заполните таблицу $D(i,c)$ для $i=0,\dots,4$ и $c=0,\dots,W$, прочитайте оптимум и пройдите назад к укладке. Укажите $D(2,W)$ и $D(4,W)$.

\textbf{Решение.}
\begingroup\sloppy
\begin{enumerate}[label=\textbf{\arabic*.}, leftmargin=*]
  \item Строка $0$ — одни нули. Строка $1$ подскакивает до $p_1$ при $c=w_1$; строка $2$ добавляет предмет $2$, но предметы $1$ и $2$ вместе слишком тяжелы, поэтому $D(2,W)=13$.
  \item Строки $3$ и $4$ по тому же правилу дают всю таблицу:
\handoutdisplay{
\begin{array}{c|ccccccccc} i\backslash c & 0 & 1 & 2 & 3 & 4 & 5 & 6 & 7 & 8 \\ \hline 0 & 0 & 0 & 0 & 0 & 0 & 0 & 0 & 0 & 0 \\ 1 & 0 & 0 & 0 & 0 & 0 & 13 & 13 & 13 & 13 \\ 2 & 0 & 0 & 0 & 0 & 10 & 13 & 13 & 13 & 13 \\ 3 & 0 & 0 & 0 & 6 & 10 & 13 & 13 & 16 & 19 \\ 4 & 0 & 0 & 5 & 6 & 10 & 13 & 15 & 18 & 19 \end{array}
}
 В углу $D(4,W)=19$.
  \item Обратный ход от $(4,W)$: $D(4,W)=D(3,W)$, значит, предмет $4$ не взят; $D(3,W)\ne D(2,W)$, значит, предмет $3$ взят, и $c$ уменьшается на $w_3$; предмет $2$ не взят, а предмет $1$ взят. Укладка: $(1;\ 3)$.
  \item Порядок по плотности взял бы предмет $1$, пропустил предмет $2$ и добавил предмет $4$: $18$, на единицу меньше оптимума. В таблице $(n+1)(W+1)=45$ клеток, каждая заполняется за постоянное время.
\end{enumerate}
\endgroup
\begin{answerbox}
$D(2,W)=13$, $D(4,W)=19$ с предметами $(1;\ 3)$.
\end{answerbox}

\problemhead{C2}{Неограниченные копии}{Самостоятельно}
\textbf{Условие.} Предметы трёх видов, в любом количестве копий, приносят $(14;\ 17;\ 11)$ и весят $(4;\ 5;\ 3)$; вместимость $W=13$. Заполните одну строку $D(c)$ для $c=0,\dots,W$ по правилу $D(c)=\max\bigl(D(c-1),\ \max_{i:\,w_i\le c}\{p_i+D(c-w_i)\}\bigr)$, затем сравните с тем, чтобы брать самый плотный вид сколько влезет и дозаполнять жадно. Укажите $D(W)$, число копий третьего вида в оптимуме и значение жадного правила.

\textbf{Решение.}
\begingroup\sloppy
\begin{enumerate}[label=\textbf{\arabic*.}, leftmargin=*]
  \item Строка равна $D=(0;\ 0;\ 0;\ 11;\ 14;\ 17;\ 22;\ 25;\ 28;\ 33;\ 36;\ 39;\ 44;\ 47)$ для $c=0,1,\dots,W$. В углу $D(W)=47$.
  \item Обратный ход по победившим предметам даёт копии $x=(1;\ 0;\ 3)$: они занимают всю вместимость.
  \item Самый плотный — третий вид. Если брать его сколько влезет, останется остаток, который ничем не заполнить, и жадная загрузка стоит $44$.
\end{enumerate}
\endgroup
\begin{answerbox}
$D(W)=47$ при $x=(1;\ 0;\ 3)$; жадное правило $44$.
\end{answerbox}

\problemhead{C3}{Точные суммы}{Самостоятельно}
\textbf{Условие.} Веса $(5;\ 6;\ 9;\ 14)$, каждый используется не больше одного раза. Отметьте, какие суммы $c=0,\dots,W$ при $W=21$ набираются точно, с помощью одной строки логических значений, которую для каждого веса обновляем от $c=W$ вниз до $c=w_i$. Укажите число достижимых сумм, включая $0$, и наибольшую из них. Равен ли оптимум рюкзака с $p_i=w_i$ числу $W$?

\textbf{Решение.}
\begingroup\sloppy
\begin{enumerate}[label=\textbf{\arabic*.}, leftmargin=*]
  \item После всех четырёх весов достижимые суммы до $W$ таковы: $(0;\ 5;\ 6;\ 9;\ 11;\ 14;\ 15;\ 19;\ 20)$, их $9$.
  \item Наибольшая — $20$, меньше $W$, поэтому ни одно подмножество не даёт в сумме $W$, и у рюкзака с $p_i=w_i$ оптимум $20$, а не $W$.
  \item В строке $W+1$ логических значений, и каждый вес проходит её один раз: снова $O(nW)$, поэтому задача о сумме подмножества, NP-полная, проста при небольшом $W$.
\end{enumerate}
\endgroup
\begin{answerbox}
Достижимых сумм $9$, наибольшая $20$; оптимум не равен $W$.
\end{answerbox}

\begin{trapbox}
$O(nW)$ — не полиномиальное время: $W$ записывается примерно $\log_2W$ битами, поэтому таблица экспоненциальна по длине входа. Если умножить все веса и вместимость на тысячу, задача не изменится, а таблица вырастет примерно в тысячу раз.
\end{trapbox}

\begin{selfcheckbox}
Когда стоит выбрать ветви и границы, а не таблицу?
\selfcheckanswer Когда веса велики или не целые, когда оценки сильные или когда важна только одна вместимость. Таблица выигрывает при небольших целых вместимостях или когда нужны все вместимости сразу.
\end{selfcheckbox}

% =====================================================================
\section{Двоичные переменные в моделях}\label{sec:modelling}

Большинство практических целочисленных моделей — линейные задачи с несколькими двоичными переключателями. Переключатели несут логику, фиксированные затраты и взаимодействия; искусство в том, чтобы записать их так, чтобы релаксация оставалась близкой к целым.

\begin{itemize}
  \item Если $A$, то $B$: $y_A\le y_B$. Не оба: $y_A+y_B\le1$. Не больше $k$ из набора: $\sum y_i\le k$. Ровно один: $\sum y_i=1$.
  \item Фиксированные затраты: $x\le u\,y$ с верной границей $u$ и $-f\,y$ в цели. Чем меньше верная граница $u$, тем сильнее релаксация.
  \item Произведение двоичных $z=x_ix_j$: $z\le x_i$, $z\le x_j$, $z\ge x_i+x_j-1$, $z\ge0$; оставляйте только ту сторону, против которой давит цель.
  \item Двоичная на непрерывную $z=y\,x$ при $0\le x\le u$: $z\le u\,y$, $z\le x$, $z\ge x-u(1-y)$, $z\ge0$.
\end{itemize}

\problemhead{D1}{Три линии и два правила}{У доски}
\textbf{Условие.} Три линии выпускают продукты с прибылью на единицу $(7;\ 8;\ 6)$; открытие линии стоит $(40;\ 20;\ 40)$. Единица продукта требует $(3;\ 4;\ 1)$ часов общей бригады, у которой $80$ часов, а спрос ограничивает выпуск значениями $(15;\ 10;\ 25)$. Линия 1 может открыться, только если открыта линия 2, и открыть можно не больше двух линий. Запишите модель с $y_j\in\{0,1\}$ и $x_j\ge0$ и решите её перебором допустимых $y$. Укажите лучшую прибыль без двух правил и с ними.

\textbf{Решение.}
\begingroup\sloppy
\begin{enumerate}[label=\textbf{\arabic*.}, leftmargin=*]
  \item Модель: $\max\sum_jp_jx_j-\sum_jf_jy_j$ при $\sum_ja_jx_j\le H$, $x_j\le u_jy_j$, $y_1\le y_2$, $y_1+y_2+y_3\le2$, $x\ge0$, $y\in\{0,1\}^3$. Строка $x_j\le u_jy_j$ сразу связывает выпуск с линией и со спросом.
  \item Без правил лучше всего линии $1$ и $3$: выпуск $(15;\ 0;\ 25)$ и прибыль $175$. Этот план нарушает первое правило.
  \item Шаблоны, подчиняющиеся обоим правилам, с лучшим выпуском для каждого:
\handoutdisplay{
\begin{array}{ccc|ccc|c} y_1 & y_2 & y_3 & x_1 & x_2 & x_3 & \text{profit} \\ \hline 0 & 0 & 0 & 0 & 0 & 0 & 0 \\ 0 & 0 & 1 & 0 & 0 & 25 & 110 \\ 0 & 1 & 0 & 0 & 10 & 0 & 60 \\ 0 & 1 & 1 & 0 & 10 & 25 & 170 \\ 1 & 1 & 0 & 15 & 8{,}75 & 0 & 115 \end{array}
}
 Лучший — линии $2$ и $3$, $x=(0;\ 10;\ 25)$, прибыль $170$.
\end{enumerate}
\endgroup
\begin{answerbox}
Без правил $175$, с ними $170$ (линии $2$ и $3$).
\end{answerbox}

\problemhead{D2}{Соседние стенды}{Самостоятельно}
\textbf{Условие.} У четырёх стендов ценности $(10;\ 4;\ 5;\ 5)$ и стоимости $(2;\ 2;\ 2;\ 5)$, бюджет $B=11$. Соседи влияют друг на друга: пары $(1,2)$, $(2,3)$ и $(3,4)$ добавляют к ценности $(3;\ -8;\ 3)$, если выбраны оба стенда пары. Замените каждое произведение $x_ix_{i+1}$ новой переменной только с нужными ей строками и решите задачу. Укажите оптимум, значение, которое выдаст модель, если для пары $(2,3)$ оставить только $z\le x_2$ и $z\le x_3$, и число использованных строк линеаризации.

\textbf{Решение.}
\begingroup\sloppy
\begin{enumerate}[label=\textbf{\arabic*.}, leftmargin=*]
  \item Двум бонусам нужны $z\le x_i$ и $z\le x_{i+1}$; цель толкает их вверх, поэтому они равны $1$ ровно тогда, когда выбраны оба стенда. Штрафу нужно $z_{2,3}\ge x_2+x_3-1$; цель толкает его вниз к $0$, если только оба не выбраны. Итого $5$ строк.
  \item Перебирая подмножества в рамках бюджета, находим лучшее: $(1;\ 3;\ 4)$ с ценностью $23$; стенд $2$ исключён, чтобы избежать штрафа.
  \item Если оставить только $z_{2,3}\le x_2$, $z_{2,3}\le x_3$, солвер может положить $z_{2,3}=0$ даже при обоих выбранных стендах, и штраф никогда не платится. Он выберет $(1;\ 2;\ 3;\ 4)$ и сообщит $30$ — значение, недостижимое ни одним реальным планом.
\end{enumerate}
\endgroup
\begin{answerbox}
Оптимум $23$ со стендами $(1;\ 3;\ 4)$; с неверными строками $30$; строк $5$.
\end{answerbox}

\begin{trapbox}
Пропущенная половина линеаризации не вызывает ошибки. Солвер просто пользуется лазейкой и сообщает значение, которого не достигает ни один реальный план, как в D2.
\end{trapbox}

% =====================================================================
\section{Сильные и слабые модели}\label{sec:formulations}

\begin{definitionbox}{Определение: Более сильная модель.}
Две модели одного и того же множества целых точек сравнивают по релаксациям: $P_1$ сильнее $P_2$, если $P_1\subseteq P_2$. Идеальная модель — выпуклая оболочка целых точек, у которой все вершины целые.
\end{definitionbox}

\begin{theorembox}{Утверждение: Оценки вложенных релаксаций.}
Если $P_1\subseteq P_2$, то для максимизации $\max_{P_1}c^{\top}x\le\max_{P_2}c^{\top}x$: более сильная модель даёт не худшую оценку, и метод ветвей и границ закрывает с ней больше узлов.
\end{theorembox}

\problemhead{D3}{Две модели одной задачи о складах}{Самостоятельно}
\textbf{Условие.} Два склада с затратами на открытие $(6;\ 8)$ обслуживают четыре магазина; обслуживание магазина $i$ со склада $j$ стоит $c_{ij}$, $C=\begin{pmatrix} 7 & 2 \\ 6 & 5 \\ 2 & 8 \\ 2 & 4 \end{pmatrix}$. Каждый магазин обслуживает один открытый склад. Запишите модель с суммарными строками $\sum_ix_{ij}\le4y_j$ и со строками $x_{ij}\le y_j$, решите обе релаксации и целочисленную задачу. Укажите оба значения ЛП и оптимум.

\textbf{Решение.}
\begingroup\sloppy
\begin{enumerate}[label=\textbf{\arabic*.}, leftmargin=*]
  \item Суммарные строки: каждый магазин идёт на свой более дешёвый склад, а $y=(0{,}5;\ 0{,}5)$ лишь покрывает их, оплачивая долю каждого открытия. Значение ЛП $18$.
  \item Раздельные строки: магазин, обслуживаемый наполовину со склада, заставляет открыть этот склад наполовину, и дробные $y$ больше ничего не экономят. Значение ЛП $23$, и его оптимум уже целый.
  \item Целый оптимум открывает только склад $(1)$ и стоит $23$ — столько же, сколько сильная оценка.
\end{enumerate}
\endgroup
\begin{answerbox}
Значения ЛП $18$ и $23$, оптимум $23$.
\end{answerbox}

\begin{whybox}
Суммарная строка $\sum_ix_{ij}\le m\,y_j$ позволяет складу полностью обслужить часть магазинов, заплатив лишь долю стоимости открытия; строки $x_{ij}\le y_j$ говорят то, что и так следует из целочисленности: магазин может пользоваться только открытым складом.
\end{whybox}

% =====================================================================
\section{Котосудоку}\label{sec:cats}

Головоломка блока G — задача допустимости в двоичных переменных: по коту на строку, столбец и область, никакие два не касаются. Решатель на странице — поиск в глубину по строкам, закрывающий узел, как только в строке или области не остаётся клеток, то есть метод ветвей и границ, где единственная оценка — недопустимость.

Как целочисленная задача она имеет строки модели из блока G. Неравенства по квадратам — более сильный выбор: на небольших полях солверу с ними обычно вовсе не нужно ветвиться.

\end{document}
